% =====================================================================
% 02-background.tex
% =====================================================================
\section{Background: From LLM RL to Agentic RL}\label{sec:background}

\subsection{The LLM Post-training Pipeline}

An LLM is first \emph{pre-trained} by maximum-likelihood next-token
prediction on large corpora, yielding a general-purpose conditional
generator. Two post-training stages then shape its behavior. In
SFT the model imitates curated
$(\text{prompt},\text{response})$ demonstrations; this is stable but
bounded by the quality of the demonstrations and offers no mechanism to
exceed them. Reinforcement Fine-Tuning (RFT) instead optimizes
the model against a reward signal, allowing it to explore the response
space and, in principle, surpass its demonstrators~\cite{ouyang2022instructgpt}.

The canonical RLHF pipeline of InstructGPT chains the two: SFT, then
training a reward model on pairwise human preferences, then policy
optimization with PPO under a Kullback--Leibler (KL) penalty that keeps
the policy close to the SFT reference~\cite{ouyang2022instructgpt,ziegler2019finetuning,stiennon2020summarize}. Notably, a 1.3B-parameter
model aligned this way was preferred by humans over the 175B GPT-3,
establishing that alignment via RL can be more effective than scaling
parameters alone.

\subsection{A Unified MDP/POMDP View}

Following Zhang \textit{et~al.}~\cite{zhang2026landscape}, we cast both
classical and agentic RL for LLMs as Partially Observable Markov Decision
Processes (POMDPs), which exposes their essential difference. An MDP is a
tuple $\langle \mathcal{S}, \mathcal{A}, P, R, T, \gamma \rangle$ with
state space $\mathcal{S}$, action space $\mathcal{A}$, transition kernel
$P$, reward $R$, horizon $T$, and discount $\gamma$.

\paragraph{PBRFT.} Classical RLHF/DPO is a
\emph{degenerate, single-step} MDP. The episode starts from one prompt
state $s_0$ and terminates immediately after the model emits a complete
response $a$:
\begin{equation}
\langle \mathcal{S}_{\text{trad}}, \mathcal{A}_{\text{trad}},
P_{\text{trad}}, R_{\text{trad}}, T{=}1, \gamma{=}1 \rangle ,
\end{equation}
with $\mathcal{S}_{\text{trad}}=\{\text{prompt}\}$, a deterministic
transition to the terminal state, and a single scalar reward
$R_{\text{trad}}(s_0,a)=r(a)$. The objective is simply
$J(\theta)=\mathbb{E}_{a\sim\pi_\theta}[r(a)]$.

\paragraph{Agentic RL.} When the LLM acts over many steps, reading tool
outputs, observing an environment, and updating memory, the process becomes
a temporally extended POMDP
$\langle \mathcal{S}_{\text{agent}}, \mathcal{A}_{\text{agent}},
P_{\text{agent}}, R_{\text{agent}}, \gamma, \mathcal{O} \rangle$ with
horizon $T>1$. The agent receives observations $o_t=\mathcal{O}(s_t)$,
chooses actions $a_t$, and the state evolves stochastically
$s_{t+1}\sim P(s_{t+1}\mid s_t,a_t)$. The action space splits into
language and environment actions,
$\mathcal{A}_{\text{agent}}=\mathcal{A}_{\text{text}}\cup\mathcal{A}_{\text{action}}$, where structured actions (delimited by
special tokens) invoke tools or modify the environment. Rewards become
step-wise, mixing sparse task completion with dense sub-rewards, and the
objective is the discounted return
\begin{equation}
J(\theta)=\mathbb{E}_{\tau\sim\pi_\theta}\Big[\textstyle\sum_{t=0}^{T-1}
\gamma^t R_{\text{agent}}(s_t,a_t)\Big] .
\end{equation}

Table~\ref{tab:pbrft-vs-agentic} summarizes the contrast. The shift from
$T{=}1$ to $T{>}1$ is the single most important conceptual move in the
field: it introduces planning, long-horizon credit assignment, and
exploration as first-class concerns.

\begin{table}[t]
\centering
\caption{Formal comparison of PBRFT and Agentic RL.}
\label{tab:pbrft-vs-agentic}
\renewcommand{\arraystretch}{1.25}
\begin{tabular}{@{}p{0.10\columnwidth}p{0.36\columnwidth}p{0.40\columnwidth}@{}}
\toprule
& PBRFT & Agentic RL \\
\midrule
State & $\{s_0\}$, single prompt; ends at once & $s_t\in\mathcal{S}$, $o_t=\mathcal{O}(s_t)$, $T>1$ \\
Action & Text sequence & $\mathcal{A}_{\text{text}}\cup\mathcal{A}_{\text{action}}$ \\
Transition & Deterministic, terminal & $P(s_{t+1}\mid s_t,a_t)$ \\
Reward & Scalar $r(a)$ & Step-wise; sparse + dense \\
Objective & $\mathbb{E}[r(a)]$ & $\mathbb{E}[\sum_t \gamma^t R(s_t,a_t)]$ \\
\bottomrule
\end{tabular}
\end{table}

\subsection{The Algorithmic Spectrum}\label{sec:background-algos}

Four algorithm families recur throughout the survey; we introduce them
here and analyze their trade-offs in Section~\ref{sec:optimization}.

\paragraph{REINFORCE.} The foundational policy
gradient~\cite{williams1992reinforce,sutton2018rl} increases the
log-probability of actions weighted by their baseline-subtracted
return,
$\nabla_\theta J=\mathbb{E}[(R(s_0,a)-b(s_0))\nabla_\theta\log\pi_\theta(a\mid s_0)]$.
It is simple but suffers from high gradient variance.

\paragraph{PPO.} PPO~\cite{schulman2017ppo}
stabilizes policy-gradient updates with a clipped surrogate objective and
a learned value function (critic) for advantage estimation
$A(s_t,a_t)=R(s_t,a_t)-V(s_t)$. PPO became the workhorse of RLHF but
requires a critic of comparable size to the policy.

\paragraph{DPO.} DPO~\cite{rafailov2023dpo}
removes the explicit reward model entirely, reformulating the
KL-constrained reward-maximization as a binary classification loss over
preferred/dispreferred pairs $(y_w,y_l)$, with the policy itself acting
as an implicit reward model.

\paragraph{GRPO.} GRPO~\cite{shao2024deepseekmath}
removes the critic: it samples a \emph{group} of responses per prompt and
normalizes their rewards within the group to form advantages. Combined
with verifiable rewards, GRPO drove the emergent reasoning of
DeepSeek-R1~\cite{guo2025deepseekr1}.
